\documentclass[10pt,a4paper]{amsart}
\usepackage{amsmath,amssymb,amsthm,graphicx}
% Retained preamble for the Stacks Project chapter snapshot in this directory.
% The chapter's own preamble.tex is a copy of the shared Stacks Project preamble; the block
% below is sliced MECHANICALLY, line for line, from it. Nothing was added, removed or
% redefined; the wrapper class, the xr-hyper/hyperref wiring and the externaldocument list
% are omitted because the assembler supplies the document class.
%
% lines taken from preamble.tex: 3, 5, 9, 11, 12, 156, 157, 158, ... (51 lines total)
%
% package coverage restored from the chapter's own preamble (2026-10-13): xy, xy-2cell

\usepackage[all]{xy}
\xyoption{2cell}
\UseAllTwocells

\newenvironment{reference}{\comment}{\endcomment}
\newenvironment{slogan}{\comment}{\endcomment}
\newenvironment{history}{\comment}{\endcomment}
\theoremstyle{plain}
\newtheorem{theorem}[subsection]{Theorem}
\newtheorem{proposition}[subsection]{Proposition}
\newtheorem{lemma}[subsection]{Lemma}
\theoremstyle{definition}
\newtheorem{definition}[subsection]{Definition}
\newtheorem{example}[subsection]{Example}
\newtheorem{exercise}[subsection]{Exercise}
\newtheorem{situation}[subsection]{Situation}
\theoremstyle{remark}
\newtheorem{remark}[subsection]{Remark}
\newtheorem{remarks}[subsection]{Remarks}
\numberwithin{equation}{subsection}
\def\lim{\mathop{\mathrm{lim}}\nolimits}
\def\colim{\mathop{\mathrm{colim}}\nolimits}
\def\Spec{\mathop{\mathrm{Spec}}}
\def\Hom{\mathop{\mathrm{Hom}}\nolimits}
\def\Ext{\mathop{\mathrm{Ext}}\nolimits}
\def\SheafHom{\mathop{\mathcal{H}\!\mathit{om}}\nolimits}
\def\SheafExt{\mathop{\mathcal{E}\!\mathit{xt}}\nolimits}
\def\Sch{\mathit{Sch}}
\def\Mor{\mathop{\mathrm{Mor}}\nolimits}
\def\Ob{\mathop{\mathrm{Ob}}\nolimits}
\def\Sh{\mathop{\mathit{Sh}}\nolimits}
\def\NL{\mathop{N\!L}\nolimits}
\def\CH{\mathop{\mathrm{CH}}\nolimits}
\def\proetale{{pro\text{-}\acute{e}tale}}
\def\etale{{\acute{e}tale}}
\def\QCoh{\mathit{QCoh}}
\def\Ker{\mathop{\mathrm{Ker}}}
\def\Im{\mathop{\mathrm{Im}}}
\def\Coker{\mathop{\mathrm{Coker}}}
\def\Coim{\mathop{\mathrm{Coim}}}
\DeclareMathSymbol{\boxtimes}{\mathbin}{AMSa}{"02}
\def\QCohstack{\mathcal{QC}\!\mathit{oh}}
\def\Cohstack{\mathcal{C}\!\mathit{oh}}
\def\Spacesstack{\mathcal{S}\!\mathit{paces}}
\def\Quotfunctor{\mathrm{Quot}}
\def\Hilbfunctor{\mathrm{Hilb}}
\def\Curvesstack{\mathcal{C}\!\mathit{urves}}
\def\Polarizedstack{\mathcal{P}\!\mathit{olarized}}
\def\Complexesstack{\mathcal{C}\!\mathit{omplexes}}
\def\Pic{\mathop{\mathrm{Pic}}\nolimits}
\def\Picardstack{\mathcal{P}\!\mathit{ic}}
\def\Picardfunctor{\mathrm{Pic}}
\def\Deformationcategory{\mathcal{D}\!\mathit{ef}}
\usepackage{cleveref}
\crefname{originallemma}{original lemma}{original lemmas}
\Crefname{originallemma}{Original lemma}{Original lemmas}
\makeatletter
\newlabel{categories-example-pro-morphism-inverse-systems}{{0G2W}{source}{}{sourcepointer-1}{}}
\newlabel{categories-example-pro-morphism-inverse-systems@cref}{{[originallemma][0][]0G2W}{[1][0][]source}{}{}{}}
\makeatother
\title{Stacks, Differential Graded Sheaves: lemma-pro-isomorphic}
\author{The Stacks Project Authors}
\date{stacks-project:sdga.tex:local-snapshot-20261012; GNU FDL 1.2 or later}
\begin{document}
\maketitle
\noindent\textit{Editorial note.} Statement span is lines 6451-6459 of sources/stacks-sdga/pinned-original-source.tex; proof span is lines 6461-6504, the complete author proof (42 body lines).\par\medskip
\noindent\textit{Editorial note.} Navigation stubs: 1 cited-but-not-included labels are supplied as non-load-bearing pointers with their pinned tags where the snapshot tag index has them.\par\medskip
\noindent\textit{Editorial note.} extra\_packages is empty on purpose: the retained preamble already loads xy with options and a second declaration clashes in the assembler wrapper.\par\medskip
\noindent\textit{Editorial note.} Text is retained verbatim, including the author's own typos and placeholders.\par\medskip
\noindent\textit{Editorial note.} Queue index 21 of queue\_sdga.json (pb=1330, stmt 6451-6459, proof 6461-6504).\par\medskip
\noindent\textit{Allowed standard prerequisites (editorial).}\begin{itemize}
\item Explicit prerequisite boundary. The cited inputs are prerequisites; the argument in the retained span is not replaced by any cited result.
\end{itemize}\medskip
\ifdefined\hypertarget\hypertarget{sourcepointer-1}{}\fi
\noindent\textit{Source reference.} Original lemma 0G2W: sources/stacks-sdga/pinned-original-source.tex, source lines 6504--6504. Cited by the retained statement or proof but not part of the collected spans; supplied only as a navigation pointer. The argument does not depend on its content.\par\medskip
\begin{lemma}
\label{lemma-pro-isomorphic}
In the situation above, suppose that $\mathcal{A} = (A_n)$ and
$\mathcal{B} = (B_n)$ are inverse systems of differential
graded $R$-algebras. If $\varphi : (A_n) \to (B_n)$ is an isomorphism
of pro-objects, then the functor $\mathit{QC}(\mathcal{A}) \to
\mathit{QC}(\mathcal{B})$ constructed above is an
equivalence.
\end{lemma}

\begin{proof}
Let $\psi : (B_n) \to (A_n)$ be a morphism of pro-objects which
is inverse to $\varphi$. According to the discussion in
Categories, Example
\ref{categories-example-pro-morphism-inverse-systems}
we may assume that $\varphi$ is given by a system of maps
as above and $\psi$ is given $n(1) < n(2) < \ldots$ and
a commutative diagram
$$
\xymatrix{
\ldots \ar[r] &
B_{n(3)} \ar[d]^{\psi_3} \ar[r] &
B_{n(2)} \ar[d]^{\psi_2} \ar[r] &
B_{n(1)} \ar[d]^{\psi_1} \\
\ldots \ar[r] &
A_3 \ar[r] &
A_2 \ar[r] &
A_1
}
$$
of differential graded $R$-algebras. Since
$\varphi \circ \psi = \text{id}$ we may, after possibly
increasing the values of the functions $n(\cdot)$ and
$m(\cdot)$ assume that $B_{n(m(i))} \to A_{m(i)} \to B_i$
is the transition map. It follows that the composition of the
functors
$$
\mathit{QC}(\mathcal{B}) \to
\mathit{QC}(\mathcal{A}) \to
\mathit{QC}(\mathcal{B})
$$
sends a good sheaf of
differential graded $\mathcal{B}$-modules $\mathcal{N} = (N_n)$
to the inverse system $\mathcal{N}' = (N'_i)$ with values
$$
N'_i = N_{n(m(i))} \otimes_{B_{n(m(i))}} B_i
$$
which is canonically quasi-isomorphic to $\mathcal{N}$
exactly because $\mathcal{N}$ is an object of
$\mathit{QC}(\mathcal{B})$
and because $N_j$ is a K-flat differential graded
module for all $j$. Since the same is true for the
composition the other way around we conclude.
\end{proof}

\end{document}
